\chapter{Introduction}
\label{ch:intro}

\section{Contexte et motivation}

La gestion des stocks reste un besoin récurrent dans les PME, ateliers et entrepôts de taille modeste : suivre les quantités disponibles, enregistrer les entrées et sorties, anticiper les ruptures et disposer d'une vue consolidée pour la décision. Les tableurs partagés atteignent rapidement leurs limites (concurrence d'édition, absence d'API, traçabilité faible, déploiement non industrialisé).

\stockflow{} répond à ce besoin par une \textbf{application web full-stack} couplée à une \textbf{chaîne DevOps reproductible} : le même dépôt source permet de développer localement (Docker Compose + Adminer), de valider automatiquement le code (GitHub Actions, incluant un déploiement éphémère sur \textbf{minikube}) et de démontrer l'orchestration Kubernetes (Helm, Prometheus, Grafana, dashboard) via \texttt{make cluster}.

\section{Périmètre du projet}

Le périmètre fonctionnel couvre :
\begin{itemize}
  \item authentification par jeton JWT avec rôles \texttt{admin} et \texttt{operateur} ;
  \item gestion des catégories et des produits (SKU, quantité, seuil d'alerte) ;
  \item mouvements de stock (\texttt{entree} / \texttt{sortie}) avec contrôle de stock négatif ;
  \item tableau de bord synthétique et liste des produits sous le seuil.
\end{itemize}

Le périmètre technique couvre :
\begin{itemize}
  \item API REST FastAPI et interface React (Vite, TypeScript, Tailwind) ;
  \item persistance PostgreSQL~16 ;
  \item conteneurisation Docker et orchestration Docker Compose pour le développement (Adminer pour la BDD) ;
  \item validation Kubernetes via \textbf{minikube} local et job CI (\texttt{values-minikube.yaml}) ;
  \item chart Helm \texttt{pfa-stock} et Terraform (validate CI ; apply prod optionnel) ;
  \item stack \textbf{kube-prometheus-stack} avec dashboard Grafana versionné ;
  \item pipelines GitHub Actions : CI (tests, minikube, validation IaC).
\end{itemize}

\section{Organisation du dépôt}

Le dépôt suit une séparation nette des responsabilités (tableau~\ref{tab:structure-depot}).

\begin{table}[htbp]
  \centering
  \caption{Structure logique du dépôt}
  \label{tab:structure-depot}
  \begin{tabular}{@{}lp{9.5cm}@{}}
    \toprule
    \textbf{Répertoire} & \textbf{Rôle} \\
    \midrule
    \texttt{backend/} & API FastAPI, modèles SQLAlchemy, Alembic, tests pytest \\
    \texttt{frontend/} & SPA React, build Vite, tests Vitest \\
    \texttt{docker/} & Compose développement local \\
    \texttt{helm/pfa-stock/} & Chart Kubernetes de l'application \\
    \texttt{infrastructure/terraform/} & Plateforme K8s (ingress, TLS, monitoring) \\
    \texttt{monitoring/} & Valeurs Prometheus/Grafana et dashboard JSON \\
    \texttt{scripts/} & Automatisation démarrage et vérification \\
    \texttt{.github/workflows/} & CI et déploiement production \\
    \texttt{docs/demo-soutenance.md} & Guide validation soutenance (Compose + minikube) \\
    \texttt{docs/vps-setup.md} & Production VPS optionnelle (annexe) \\
    \bottomrule
  \end{tabular}
\end{table}

\section{Modes d'exécution}

Deux modes coexistent :

\begin{description}
  \item[\texttt{make dev}] Stack Docker Compose : frontend sur le port~3000, API sur le port~8000, Adminer sur le port~8080.
  \item[\texttt{make cluster}] Cluster \textbf{minikube} local : chart Helm, kube-prometheus-stack, dashboard Kubernetes (voir \texttt{docs/demo-soutenance.md}).
  \item[CI GitHub Actions] Six jobs dont \texttt{kubernetes} (minikube éphémère + smoke test \texttt{/health}).
\end{description}

Un déploiement production sur VPS (\texttt{deploy.yml}, \texttt{docs/vps-setup.md}) reste documenté comme \textbf{perspective optionnelle}, hors périmètre de validation devant le jury.

\section{Organisation du rapport}

Le chapitre~\ref{ch:besoin} formalise les objectifs et les acteurs. Le chapitre~\ref{ch:archi} présente l'architecture globale et les flux principaux. Les chapitres~\ref{ch:backend} à~\ref{ch:donnees} détaillent l'application ; les chapitres~\ref{ch:docker} à~\ref{ch:terraform} l'infrastructure ; les chapitres~\ref{ch:monitoring} et~\ref{ch:cicd} l'observabilité et l'intégration continue. Le chapitre~\ref{ch:exploitation} documente l'exploitation quotidienne. Les annexes recensent l'API REST et les variables de configuration.

\section{Conventions documentées}

\begin{itemize}
  \item Les chemins de fichiers sont relatifs à la racine du dépôt sauf mention contraire.
  \item Les identifiants de démonstration (\texttt{admin@stock.tn}, etc.) sont documentés à titre opérationnel ; en production ils doivent être remplacés et stockés dans un gestionnaire de secrets (GitHub Secrets, Secrets Kubernetes).
\end{itemize}
